\documentclass[11pt]{article}
\usepackage[T1]{fontenc}
\usepackage{lmodern}
\usepackage[margin=0.95in]{geometry}
\usepackage{amsmath,amssymb}
\usepackage{microtype}
\usepackage{enumitem}
\usepackage{xcolor}
\usepackage[colorlinks=true,linkcolor=blue!50!black]{hyperref}

\newcommand{\C}{\mathbb{C}}
\newcommand{\R}{\mathbb{R}}
\newcommand{\T}{\mathbb{T}}
\newcommand{\inner}[2]{\langle #1,#2\rangle}
\newcommand{\abs}[1]{\lvert #1\rvert}
\newcommand{\norm}[1]{\lVert #1\rVert}
\newcommand{\key}[1]{\par\medskip\noindent{\normalsize\bfseries #1}\par\nobreak\vspace{-1pt}}
\setlist[enumerate]{leftmargin=1.5em,itemsep=3pt,topsep=3pt}

\title{\bfseries Step 1 of 5: from a continuum to finitely many balls\\[5pt]
\large\normalfont Companion to \texttt{mub-outline.pdf}, Figure 1;
preprint \S2 (Lemmas 2.2--2.9, Proposition 2.10)}
\author{}
\date{}

\begin{document}
\maketitle
\vspace{-2.6em}

\noindent
\textbf{What this step receives.} Every dephased complex Hadamard matrix $H$ --- $25$ free phases
modulo the symmetries of Lemma 2.2 --- together with no hypothesis on $H$ whatsoever.

\par\medskip\noindent\textbf{What it must produce.} A \emph{finite} list of centres $u^{(1)},\dots,u^{(m)}$ such
that every exact $H$ has its nine seed phases within $R_s=2\pi(0.0320001)$ radians of one of them, and
such that the list is short enough to visit.

\key{1.\ Spend the symmetries before cutting anything}

Eleven of the $36$ phases are pure gauge: permute rows and columns and multiply them by unit scalars
until the first row and first column are $1$. Two more symmetries matter later --- transposition and
entrywise conjugation --- and one of them is \emph{needed} to make the next step work. The point is
economy: the smaller the quotient, the fewer boxes.

\key{2.\ Arrange canonically, and only then bin}

Lemma 2.4 is the step that makes the search finite, and it is the one worth teaching. Among all pairs
of rows and all pairs of columns, take a pair whose six phase differences have the smallest
\emph{gap} (the largest circular arc containing none of the six points). Transpose if it is a column
pair, permute rows so it is the first two, order the six phases around the circle starting just after
a largest gap, and reverse the direction if the orientation condition $x_1\le x_5-x_4$ fails. The
result is
\[
0\le x_1\le\dots\le x_5\le 1,\qquad g_x=1-x_5,\qquad x_1\le x_5-x_4 ,
\]
and similarly for a second column, sharing one entry $\xi$ with that row. Now a box in the nine seed
phases is a \emph{multiset} of bin indices, so the count collapses:
\[
112^9=2.8\times10^{18}\quad\longrightarrow\quad\binom{116}{5}=1.6\times10^{8}.
\]
Canonical form first, bins second: the other order is unaffordable by ten orders of magnitude.

\key{3.\ Bin the circle, and remember what a box is}

Bins $I_b=[b/N,(b{+}1)/N]$ with $N=112$ are \emph{closed}, so the cover misses no phase; a boundary
phase belongs to two boxes, which is harmless. A box is a product of intervals --- the only kind of
region for which estimates are cheap --- and its midpoint $u$ is a rational point that need not be a
Hadamard matrix. The circumradius of a nine-dimensional box of side $w$ is $\tfrac32w=0.0134$, which is
why a box can fit inside a ball of radius $0.032$.

\key{4.\ Four necessary tests, each with a margin}

A test may discard a box only if every exact matrix in that box fails it. That is why each test is
proved for the \emph{centre} and then widened by the box half-width:
\begin{enumerate}[label=(\alph*)]
\item \emph{Gaps.} Lemma 2.5 says a gap computed from phases is Lipschitz in those phases with a
constant of order $2$; so a gap comparison true at the centre with room to spare is true on the box.
Lemma 2.8's direction conditions sharpen this into $g_{x,\mathrm{mid}}\le1-\frac{l-1}{N}$ and
$g_{y,\mathrm{mid}}\ge1-\frac{l+2}{N}$, integer inequalities that most boxes fail.
\item \emph{Vanishing sums.} Two rows are orthogonal exactly when $1+\sum_je(x_j)=0$. Write the
midpoint residual as a real two-vector $c$, its $2\times5$ linearisation as $J$, and use
$1+\sum_je(u_j+\Delta_j)=0$ with $\abs{\Delta_j}\le w/2$ to get
\[
\norm{c+J\Delta}\le \tfrac{5\pi w^2}{4}=3.1\times10^{-4}.
\]
A box whose data cannot be repaired by any displacement of size $w/2$ is discarded. The tolerance is
quadratic in $w$ because the linear term is exactly what $J$ accounts for.
\item \emph{Supports.} To test that no $\Delta$ works, the code proposes a vector $v$ and evaluates the
dual inequality; validity does not depend on how $v$ was found, so a heuristic solve cannot spoil it.
\item \emph{Reuse.} A box whose feasible part already lies in a previously retained ball is discarded,
certified by $\norm{k+h}^2\le\norm{k}^2+\tfrac94+2k\cdot h$ with $k=N(u-u')$ an integer bin
displacement. This is what keeps the list short; without it the cover would be correct and useless.
\end{enumerate}

\key{5.\ What comes out, and what it does not say}

Proposition 2.10: the retained centres, each with its recorded selected position, cover the seed
coordinates of every exact Hadamard matrix up to the equivalences of Lemma 2.2, by closed Euclidean
balls of radius $0.032$. Note the shape of the claim. It says nothing about the lower-right block
$D'$; that sixteen-phase problem is handed to Step 2 untouched. It is stated for \emph{exact}
matrices, not for approximations. And it is a statement about a run: a box that fails a check is
retained as \emph{unresolved}, never dropped, so the only way to lose a matrix is a bug or an
arithmetic failure, both of which Appendix A bounds.

\key{Where this step can fail, and why that is safe}

Every failure mode here is harmless by construction. A test too weak to reject retains boxes and the
search gets slower; a frame or rank check that fails retains a region as unresolved and the proof is
merely not finished. No branch of this step can \emph{exclude} a region that contains a matrix, because
each test is a necessary condition for existence. That asymmetry --- disjunction on the covering side,
conjunction on the excluding side --- is the whole logical content of Stage 1.

\medskip
\noindent\textit{Exercise.} Show that $\binom{116}{5}$ counts the boxes of one list, and identify which
fraction of them survives the orientation condition $a\le l-d+1$. (Exact answer: $0.53$.)

\end{document}
